\documentclass[11pt]{article}
\usepackage[margin=1in]{geometry}
\usepackage{amsmath,amssymb,amsthm}
\usepackage{graphicx}
\usepackage{hyperref}

\newtheorem{theorem}{Theorem}
\newtheorem{lemma}{Lemma}
\newtheorem{corollary}{Corollary}
\newtheorem{remark}{Remark}

\newcommand{\K}{\mathcal K}
\newcommand{\Lcal}{\mathcal L}
\newcommand{\what}{\widehat{\otimes}}

\title{\(M\)-Embedded Upgrade of Rao's Theorem on Nice Surjections}
\author{Run \texttt{fa\_banach\_001}, agent \texttt{agent\_lane\_04}}
\date{June 18, 2026}

\begin{document}
\maketitle

\section*{Source and Target}

The source paper is T. S. S. R. K. Rao, \emph{Nice surjections on spaces of operators}, arXiv:0707.2140. Remark 14 asks whether Rao's preceding theorem remains true if \(Y\) is assumed merely to be an \(M\)-embedded space with strictly convex dual. The remark identifies the missing point:

\begin{center}
\includegraphics[width=0.95\textwidth]{figures/open_problem_crop.png}
\end{center}

In text, the source says that to complete the proof one would require an into isometry of \(Y^*\) to be \(\sigma(Y^*,Y)\)-continuous.

\section*{Result}

We prove that the missing isometry statement is true and then insert it into Rao's proof. Since a nonreflexive \(M\)-embedded \(Y\) has
\[
\K(X,Y)^{**}\cong \Lcal(X,Y^{**}),
\]
the natural upgraded version of Rao's theorem is a theorem on \(\Lcal(X,Y^{**})\). When \(Y\) is reflexive this is exactly Rao's original displayed target space \(\Lcal(X,Y)\).

\section*{Idea of the Proof}

Rao's proof of Theorem 13 first restricts the adjoint of the operator on \(\K(X,Y)^{**}\) to \(\K(X,Y)^*=X\what_{\pi}Y^*\). Then Rao's Theorem 9 gives a tensor-factor representation
\[
\Psi(x\otimes y^*)=Vx\otimes Jy^*,
\]
where \(V:X\to X\) and \(J:Y^*\to Y^*\) are into isometries. In the reflexive case, \(J\) is automatically weak-star continuous and hence has a preadjoint \(U:Y\to Y\).

For \(M\)-embedded \(Y\), this automatic continuity is the only missing point. The decomposition
\[
Y^{***}=Y^*\oplus_1Y^\perp
\]
forces every into isometry \(J:Y^*\to Y^*\) to preserve the predual \(Y\): if \(J^{**}\eta\) had a nonzero \(Y^*\)-component for some \(\eta\in Y^\perp\), then the \(\ell_1\)-norm identity would force that one vector to be norm-additive with all vectors in \(J(Y^*)\), contradicting strict convexity after applying the identity to \(u\) and \(-u\). Thus \(J\) is weak-star continuous, so \(J=U^*\). The rest of Rao's proof then gives \(\Phi=\psi^{**}\), where \(\psi(T)=UTV\).

\section*{The Full Upgrade}

\begin{theorem}
Let \(X\) be reflexive and let \(Y\) be an \(M\)-embedded Banach space. Assume that \(X\) and \(Y^*\) are strictly convex, that \(X\) is not isometric to a subspace of \(Y^*\), and that \(Y^*\) is not isometric to a subspace of \(X^*\). Assume also that the canonical identifications used in Rao's proof hold:
\[
\K(X,Y)^*=X\what_{\pi}Y^*,\qquad
\K(X,Y)^{**}=\Lcal(X,Y^{**}).
\]
Suppose \(\Phi:\Lcal(X,Y^{**})\to\Lcal(X,Y^{**})\) is a bounded linear surjection such that \(\Phi^*\) maps weak-star-continuous extreme points of the dual unit ball to extreme points.

Then there are \(U\in\Lcal(Y)\) and \(V\in\Lcal(X)\), with \(U^*:Y^*\to Y^*\) and \(V:X\to X\) into isometries, such that
\[
\Phi(S)=U^{**}SV\qquad(S\in \Lcal(X,Y^{**})).
\]
Equivalently, \(\Phi\) is the bitranspose of the nice surjection
\[
\psi:\K(X,Y)\to\K(X,Y),\qquad \psi(T)=UTV.
\]
\end{theorem}

\begin{proof}
Put \(E=\K(X,Y)\). Under the stated identifications, \(E^*=X\what_{\pi}Y^*\) and \(E^{**}=\Lcal(X,Y^{**})\). Since \(Y\) is \(M\)-embedded, \(Y^*\) has the Radon--Nikodym property; together with the corresponding hypothesis on \(X\) and the metric approximation assumptions used in Rao's proof, this is the setting in which the source identifies \(E^*\) with \(X\what_{\pi}Y^*\) and uses that the unit ball of \(E^*\) is the norm-closed convex hull of its extreme points.

As in Rao's proof of Theorem 13, the hypothesis on weak-star-continuous extreme points implies that
\[
\Psi:=\Phi^*|_{E^*}:E^*\to E^*
\]
is a bounded one-to-one operator preserving extreme points of the unit ball. Rao's Theorem 9, with the same strict convexity and non-isometric-subspace assumptions, gives into isometries \(V:X\to X\) and \(J:Y^*\to Y^*\) such that
\[
\Psi(x\otimes y^*)=Vx\otimes Jy^*
\qquad(x\in X,\ y^*\in Y^*).
\]
The only place where the reflexivity of \(Y\) was used in the original proof is to regard the into isometry \(J\) as the adjoint of an operator on \(Y\). We prove that this remains true for \(M\)-embedded \(Y\) in the lemma below. Therefore there is \(U\in\Lcal(Y)\) with \(J=U^*\).

Define \(\psi:\K(X,Y)\to\K(X,Y)\) by
\[
\psi(T)=UTV.
\]
For \(x\in X\), \(y^*\in Y^*\), and \(T\in\K(X,Y)\),
\[
\langle x\otimes y^*,\psi(T)\rangle
=y^*(UTVx)
=U^*y^*(TVx)
=\langle Vx\otimes U^*y^*,T\rangle.
\]
Thus \(\psi^*=\Psi\). Since \(\Psi\) preserves extreme points, \(\psi\) is nice; since \(\Psi\) is an into isometry, \(\psi\) is a quotient map and hence a nice surjection.

Finally \(\Phi^*\) and \(\psi^{***}\) agree on the canonical copy of \(E^*\) in \(E^{***}\). Both maps are weak-star continuous on \(E^{***}\) with predual \(E^{**}\), and \(E^*\) is weak-star dense in \(E^{***}\). Hence \(\Phi^*=\psi^{***}\), so \(\Phi=\psi^{**}\). Under \(E^{**}=\Lcal(X,Y^{**})\), the bitranspose of \(T\mapsto UTV\) is
\[
S\mapsto U^{**}SV,
\]
which proves the theorem.
\end{proof}

\section*{The Missing Lemma}

\begin{lemma}
Let \(Y\) be an \(M\)-embedded Banach space and assume that \(Y^*\) is strictly convex. Then every linear isometry
\[
J:Y^*\longrightarrow Y^*
\]
is \(\sigma(Y^*,Y)\)-continuous.
\end{lemma}

\begin{proof}
The statement is trivial if \(Y^*=\{0\}\), so assume \(Y^*\neq\{0\}\).
Since \(Y\) is \(M\)-embedded,
\[
Y^{***}=Y^*\oplus_1Y^\perp,
\]
where \(Y^\perp=\{\eta\in Y^{***}:\eta(y)=0\text{ for all }y\in Y\}\).

Let \(J:Y^*\to Y^*\) be a linear isometry. Its bitranspose
\[
J^{**}:Y^{***}\to Y^{***}
\]
is again a linear isometry, and \(J^{**}\) maps the canonical copy of \(Y^*\) into itself.

We claim that \(J^{**}(Y^\perp)\subseteq Y^\perp\). Fix \(\eta\in Y^\perp\), and write
\[
J^{**}\eta=a+b,\qquad a\in Y^*,\quad b\in Y^\perp.
\]
For every \(u\in Y^*\), the \(\ell_1\)-sum norm gives
\[
\|u\|+\|\eta\|
=\|u+\eta\|
=\|J^{**}u+J^{**}\eta\|
=\|Ju+a\|+\|b\|.
\]
Also
\[
\|\eta\|=\|J^{**}\eta\|=\|a\|+\|b\|.
\]
Subtracting yields
\[
\|Ju+a\|=\|Ju\|+\|a\|\qquad\text{for all }u\in Y^*.
\]
If \(a\neq0\), choose \(u\in Y^*\) with \(Ju\neq0\). Strict convexity of \(Y^*\) forces equality in the triangle inequality only along a single nonnegative ray, so \(Ju\) and \(a\) must be nonnegative scalar multiples of one another. Applying the same equality to \(-u\) forces \(-Ju\) and \(a\) to be nonnegative scalar multiples of one another. This is impossible. Hence \(a=0\), and \(J^{**}\eta\in Y^\perp\).

Thus \(Y^\perp\) is invariant under \(J^{**}\). For \(y\in Y\) and \(\eta\in Y^\perp\),
\[
\eta(J^*y)=J^{**}\eta(y)=0.
\]
Therefore \(J^*y\in (Y^\perp)_\perp=Y\). Hence \(J^*(Y)\subseteq Y\), so \(J\) has a preadjoint \(J_*:Y\to Y\). Equivalently, \(J=(J_*)^*\), and \(J\) is \(\sigma(Y^*,Y)\)-continuous.
\end{proof}

\begin{corollary}
Rao's Remark 14 has a positive answer in the natural bidual formulation: the proof of Theorem 13 extends from reflexive \(Y\) to \(M\)-embedded \(Y\) with strictly convex dual after replacing \(\Lcal(X,Y)\) by \(\Lcal(X,Y^{**})=\K(X,Y)^{**}\). If \(Y\) is reflexive, this is exactly the original target space.
\end{corollary}

\begin{proof}
The lemma supplies the only missing weak-star continuity step named in Remark 14. The theorem above then repeats Rao's proof of Theorem 13 with that step inserted.
\end{proof}

\section*{Verification and Limitations}

The proof is formal once the tensor identifications in the theorem statement are available. The packet states these identifications explicitly because the source's approximation-property hypotheses are serving precisely this role. No computational verification is relevant.

The only interpretive caveat is target-space notation: for nonreflexive \(M\)-embedded \(Y\), the natural bidual of \(\K(X,Y)\) is \(\Lcal(X,Y^{**})\), not \(\Lcal(X,Y)\). This packet therefore proves the theorem in the natural form suggested by the paragraph immediately preceding Rao's Theorem 13.

Bounded novelty check: run indexes had no full packet for arXiv:0707.2140. Web searches on June 18, 2026 for the exact title, for the exact Remark 14 phrase, and for close phrases such as ``M-embedded strictly convex dual into isometry weak-star continuous'' did not reveal a later published answer. Local source-index searches likewise found only the original Remark 14 and this run's earlier partial packet.

\section*{Human Review Recommendation}

Review as a claimed full solution to Rao's Remark 14 in the natural bidual formulation. The highest-value checks are:
\begin{enumerate}
\item confirm the tensor identifications under the intended approximation-property hypotheses;
\item confirm that Rao's Theorem 9 proof uses reflexivity of \(Y\) only to obtain a preadjoint for the \(Y^*\)-isometry;
\item confirm the bitranspose formula \( \psi^{**}(S)=U^{**}SV \).
\end{enumerate}

\begin{thebibliography}{9}

\bibitem{Rao}
T. S. S. R. K. Rao,
\emph{Nice surjections on spaces of operators},
arXiv:0707.2140.

\bibitem{HWW}
P. Harmand, D. Werner, and W. Werner,
\emph{\(M\)-ideals in Banach spaces and Banach algebras},
Lecture Notes in Mathematics 1547, Springer, 1993.

\bibitem{DU}
J. Diestel and J. J. Uhl,
\emph{Vector Measures},
Mathematical Surveys 15, American Mathematical Society, 1977.

\end{thebibliography}

\end{document}
